\section*{Overview}

Knapsack is the template DP for resource allocation. The items change from problem to problem, but the state shape is
stable: after considering some objects, what is the best value achievable with a given amount of capacity?

\section*{When to Use It}

Use knapsack-style DP when:

\begin{itemize}
  \item choices consume some bounded resource,
  \item items can be taken zero, one, or many times,
  \item the resource limit is small enough for pseudo-polynomial DP.
\end{itemize}

\section*{Core Idea}

Let \(\texttt{dp[w]}\) be the best value achievable with total weight exactly or at most \(w\). Each item updates this
array by either using the item or skipping it.

\section*{Key Insight}

The direction of the capacity loop is the entire difference between 0/1 and unbounded knapsack:

\begin{itemize}
  \item descending loop: each item used at most once,
  \item ascending loop: the same item may contribute again in the same phase.
\end{itemize}

\section*{Operations / Main Technique}

\begin{itemize}
  \item choose the state meaning carefully,
  \item initialize impossible states correctly,
  \item update by item transitions,
  \item compress the first dimension when only the previous item layer is needed.
\end{itemize}

\paragraph{Worked example.}
In 0/1 knapsack, if an item has weight \(3\) and value \(7\), then every capacity \(w \ge 3\) can try the transition
\[
dp[w] = \max(dp[w], dp[w - 3] + 7),
\]
but only from the previous-item layer, which is why the loop goes backward.

\section*{Correctness Intuition}

Every state represents the best answer among all valid subsets respecting the processed prefix of items. The transition
simply partitions those subsets into two groups: ones that do not use the current item and ones that do.

\section*{Complexity Analysis}

With capacity \(W\) and \(n\) items, the standard implementation is \(O(nW)\) time and \(O(W)\) memory after
compression.

\section*{Implementation}

The code includes:

\begin{itemize}
  \item 0/1 knapsack,
  \item unbounded knapsack.
\end{itemize}

Those two loop directions are the main thing to remember.

\section*{Common Pitfalls}

\begin{itemize}
  \item Using the wrong loop direction.
  \item Initializing impossible states as zero when they should be negative infinity.
  \item Forcing knapsack on constraints where \(W\) is too large and another state definition is needed.
\end{itemize}

\section*{Variants / Extensions}

\begin{itemize}
  \item Bounded knapsack.
  \item Value-based knapsack.
  \item Multiple constraints or grouped items.
\end{itemize}

\section*{Practice Problems}

\begin{itemize}
  \item Standard 0/1 knapsack.
  \item Unbounded coin change or resource allocation.
  \item Value-based DP when weights are too large.
\end{itemize}

\section*{References}

\begin{itemize}
  \item \href{https://cp-algorithms.com/dynamic_programming/knapsack.html}{cp-algorithms: Knapsack Problem}.
  \item \href{https://usaco.guide/gold/knapsack}{USACO Guide: Knapsack DP}.
  \item \href{https://codeforces.com/catalog}{Codeforces Catalog}.
\end{itemize}
